% ==================================================
% Información del documento
% ==================================================

% Descomentar el idioma que desea usar en el documento, comente (%) el idioma que no desea usar
% \setLanguage{spanish}
\setLanguage{english}


\titleinfo{Título y/o subtítulo del trabajo de grado}
\titleshort{Título del trabajo de grado ...}       % Reemplaza con las primeras palabras de tu título o palabras significativas del mismo, seguido de puntos suspensivos. No permitas que título y número de página pasen a una segunda línea o más. Si no es tan extenso escriba el título completo en el campo anterior sin puntos suspensivos


 % No omitir ningún nombre o apellido; no abreviar ni dejar solo iniciales.
 % En caso de que desee agregar mas autores los debe agregar de forma manual en la linea 23 del archivo '[0.1]-Portada.tex'.
\authorone{Nombre autor 1}{Apellido autor 1}
\authortwo{}{}      % NO USADO


\worktype{(tipo de trabajo)}
\degree{(título otorgado por la UdeA)}


\advisortype{(Tipo de orientador)}
\advisor{(nombre del orientador)}
\advisordegree{(título académico mas alto)}


\faculty{(facultad, escuela, instituto o corporación UdeA)}
\program{(pregrado o posgrado UdeA)}
\city{(ciudad UdeA)}
\gradyear{(año de grado)}        % Determina con tu asesor cuál es el año para registrar: año de entrega del documento o año en que se recibe el título académico. 


\rectorName{Nombres y Apellidos}
\deanName{Nombres y Apellidos}
\departmentHead{Nombres y Apellidos}


%% Opcionales
\cohort{(cohorte posgrado en numero romano)}                % NO USADO
\researchGroup{(Grupo de investigación UdeA (A-Z))}         % NO USADO
\researchCenter{(Centro de investigación UdeA (A-Z))}       % NO USADO
\libraryName{(Biblioteca, CRAI o centro de documentación UdeA (A-Z))}